% ============================================================================
% Fichier : annexes/Ann_D_textes.tex
% Titre   : Textes Législatifs et Réglementaires
% Version : 1.0 -- novembre 2026
% ============================================================================
% Extraits identiques à ceux des chapitres (prologue, C01, C03, C04, C07,
% C10). Loi n° 2010/012 : version anglaise officielle. COBAC : relevé en
% ligne, à collationner. AUSCGIE, loi n° 2018/012, loi n° 2024/017 :
% présentés d'après des sources secondaires, à vérifier.
% ============================================================================

\chapter{Textes Législatifs et Réglementaires}
\label{ann:textes}

Cette annexe rassemble les textes cités dans le cours, avec leur degré de
vérification. Seuls les extraits lus dans leur version officielle sont
reproduits entre guillemets~; les autres sont présentés sans citation.

\section*{Degré de vérification}

\begin{tableaucours}{L{4.6cm}L{3.8cm}L{3.8cm}}{%
    \tableauHeader{Texte} &
    \tableauHeader{Vérification} &
    \tableauHeader{Chapitres}
}
Loi n°~2010/012 (cybersécurité) & Version anglaise officielle lue &
Prologue, \ref{chap:referentiels}, \ref{chap:cobit-conception},
\ref{chap:droit} \\
Règlement COBAC R-2016/04 & Art.~1 et~2 relevés en ligne~; à collationner &
\ref{chap:droit} \\
AUSCGIE révisé (2014), art.~829-1 & D'après la présentation de la réforme
par l'OHADA & \ref{chap:gouv-entreprise}, \ref{chap:droit} \\
Loi n°~2018/012 (régime financier), art.~69 & D'après la présentation du
MINFI~; à vérifier au Journal officiel & \ref{chap:alignement} \\
Loi n°~2024/017 (données personnelles) & D'après les présentations
publiées~; à vérifier au Journal officiel & \ref{chap:cobit-conception},
\ref{chap:droit} \\
Convention de Malabo (UA, 2014) & Entrée en vigueur vérifiée~; statut du
Cameroun non établi & \ref{chap:droit} \\
\end{tableaucours}

\section*{Loi n° 2010/012 relative à la cybersécurité et à la cybercriminalité}

\begin{boxjuris}
\textbf{Article~7, alinéa~2}~: \textit{« The Agency \textelp{} shall be
responsible for the regulation, control and monitoring of activities related
to the security of electronic communication networks, information systems,
and electronic certification on behalf of the State. »}

\medskip
\textbf{Article~11} (extrait)~: \textit{« The following activities may be
subject to authorization: \textelp{} the provision of security auditing,
security programs editing, and other authorized security services to the
public. »}

\medskip
\textbf{Article~26, alinéa~1}~: \textit{« Operators of information systems
shall take every technical and administrative measure to ensure the security
of services offered. To this end, they shall have standardized systems
enabling them to at all times identify, assess, process or manage any risk
relating to the security of the information systems of the services
provided directly or indirectly. »}

\medskip
\textbf{Article~29, alinéa~1}~: \textit{« Operators of information systems
shall be bound to conserve the connection and traffic data of their
information systems for a period of 10 (ten) years. »}

\medskip
\textbf{Article~30, alinéa~1}~: \textit{« Operators of information systems
shall assess and revise their security systems and, where necessary, make
the appropriate modifications to their security practices, measures and
techniques according to technological change. »}

\medskip
\textbf{Article~31, alinéa~1}~: \textit{« Electronic communication networks
and information systems content providers shall be bound to ensure the
availability of material, as well as the data stored in their
installations. »}

\medskip
\textbf{Article~32, alinéas~1 et~2}~: \textit{« (1) Electronic
communication networks and information systems shall be subject to a regime
of compulsory and periodic auditing of their security systems by the
Agency. (2) Security audit and severity scale rating shall be undertaken
each year or as required by the prevailing circumstances. »}

\medskip
\textbf{Articles~41 et~42}~: \textit{« 41. Every individual shall have the
right to the protection of their privacy. \textelp{} 42. The confidentiality
of information channelled through electronic communication and information
systems networks, including traffic data, shall be ensured by operators of
electronic communication and networks information systems. »}

\medskip
\textbf{Article~64, alinéas~1 à~3}~: \textit{« (1) Corporate bodies shall be
criminally liable for offences committed on their account by their
management structures. (2) The criminal liability of corporate bodies shall
not preclude that of natural persons who commit such offences or are
accomplices. (3) The penalties to be meted out on defaulting corporate
bodies shall be fines of from 5~000~000 (five million) to 50~000~000 (fifty
million) CFA francs. »}
\end{boxjuris}

L'article~13 soumet en outre les systèmes d'information des opérateurs à un
audit de sécurité obligatoire, dont les modalités sont fixées par voie
réglementaire~; il est présenté sans citation (Prologue).

\section*{Règlement COBAC R-2016/04}

\begin{boxjuris}
\textbf{Articles~1 et~2} (texte relevé sur une base réglementaire en ligne,
à collationner avec la version du SG-COBAC)~: article~1~: \textit{« Le
présent règlement fixe les exigences minimales relatives au système de
contrôle interne que les établissements assujettis doivent mettre en
œuvre. »}~; article~2, risque opérationnel~: \textit{« risque de pertes
résultant de l'inadéquation ou de la défaillance des processus internes, des
personnes et des systèmes, ou d'évènements externes. »}
\end{boxjuris}

\section*{Textes présentés sans citation}

\begin{itemize}
    \item \textbf{AUSCGIE révisé, article~829-1}~: selon la présentation de
          la réforme, comité d'audit obligatoire dans les sociétés cotées,
          chargé notamment du suivi de l'efficacité du contrôle interne et
          de la gestion des risques (Chapitre~\ref{chap:gouv-entreprise}).
          Responsabilité des dirigeants~: articles~161 et suivants, 740,
          d'après la doctrine (Chapitre~\ref{chap:droit}).
    \item \textbf{Loi n°~2018/012, article~69}~: selon le MINFI, rôle du
          responsable de programme dans le budget par programmes
          (Chapitre~\ref{chap:alignement}).
    \item \textbf{Loi n°~2024/017}~: protection des données à caractère
          personnel~; autorité de protection, encadrement des transferts,
          délai de mise en conformité de dix-huit mois, d'après les
          présentations publiées (Chapitres~\ref{chap:cobit-conception}
          et~\ref{chap:droit}).
    \item \textbf{Convention de Malabo}~: cybersécurité et protection des
          données, en vigueur depuis le 8~juin~2023~; ratification par le
          Cameroun non établie (Chapitre~\ref{chap:droit}).
\end{itemize}
